
Iterative LLM code repair relies primarily on execution feedback---test results and stack traces. Static analysis offers complementary signal: mechanistic explanations of why code is problematic, not just where it fails. Prior work demonstrates that static analysis improves code quality metrics (security, readability), but the impact on functional correctness (pass@k) remains unmeasured.

This paper presents a methodology study for evaluating static analysis in LLM code repair. A gate-based hypothesis validation framework is developed alongside a pylint analysis pipeline for HumanEval. The existence gate tests whether static analyzers produce actionable warnings on benchmark code---a prerequisite for any intervention to have signal.

The primary finding is that canonical HumanEval solutions exhibit only 9.15\% pylint warning rate (15 of 164 problems with at least one warning, 23 total warnings across 9 categories). This result makes canonical solutions unsuitable proxies for LLM-generated code, which typically contains more issues. The gate correctly halted downstream experiments before resources were wasted on invalid methodology.

Contributions include: (1) the first baseline measurement of static warning rates on HumanEval canonical solutions; (2) a validated pylint pipeline for code generation research; (3) a gate-based validation methodology that catches methodology limitations early; and (4) identification of the canonical-vs-LLM proxy limitation. The hypothesis that static analysis improves pass@k remains untested. The validated infrastructure and baseline are provided for future work with LLM API access.
